\documentclass[acmsmall,nonacm,screen]{acmart}

\usepackage{mathpartir}
\usepackage{booktabs}

\citestyle{acmauthoryear}
\settopmatter{printacmref=false}

\AtEndPreamble{\theoremstyle{acmdefinition}\newtheorem{remark}[theorem]{Remark}}

\newcommand{\kw}[1]{\mathsf{#1}}
\newcommand{\nm}[1]{\mathit{#1}}
\newcommand{\T}{\kw{T}}
\newcommand{\V}{\kw{V}}
\newcommand{\U}{\kw{U}}
\newcommand{\lev}[1]{\ell(#1)}
\newcommand{\ent}[2]{#1\langle #2 \rangle}
\newcommand{\op}[1]{\kw{#1}}
\newcommand{\keep}[1]{#1{\downarrow}}

% A sentence that states a fact in words. The script has one check for every name.
\newcommand{\claim}[2]{#2}

% The rows of a table that the script writes.
\makeatletter
\newcommand{\rows}[1]{\@@input artifact/expected/#1.tex }
\makeatother

% written by outlives.py; do not edit
\newcommand{\exCostChosen}{1,089}
\newcommand{\exCostLength}{12}
\newcommand{\exCostSessions}{200}
\newcommand{\exCostShare}{53}
\newcommand{\exCostSteps}{2,045}
\newcommand{\exDecayPeriods}{20,000}
\newcommand{\exLines}{600}
\newcommand{\exLookupFutures}{131,072}
\newcommand{\exMutants}{five}
\newcommand{\exOutputs}{three}
\newcommand{\exPlainRuns}{133,822}
\newcommand{\exPrograms}{3,000}
\newcommand{\exProgramsAll}{6,000}
\newcommand{\exRetryFive}{5}
\newcommand{\exRetryTrials}{4,000}
\newcommand{\exRetryWhole}{99}
\newcommand{\exSeed}{2026}
\newcommand{\exSessionActs}{24}
\newcommand{\exSessionAfter}{456}
\newcommand{\exSessionAgents}{five}
\newcommand{\exSessionCalls}{504}
\newcommand{\exSessionFirst}{seventh}
\newcommand{\exSessionLookups}{17}
\newcommand{\exSessionPages}{20}
\newcommand{\exSessionReads}{42}
\newcommand{\exSessionShare}{90}
\newcommand{\exSessionSummaries}{three}
\newcommand{\exShortOps}{nine}
\newcommand{\exShortPrograms}{7,380}
\newcommand{\exShortRuns}{47,288}
\newcommand{\exTryLeast}{18}
\newcommand{\exTryMost}{181}
\newcommand{\exTryRelabel}{64}
\newcommand{\exTryTypedText}{181}
\newcommand{\exTryWashRecord}{64}
\newcommand{\exTryWashStore}{21}
\newcommand{\exTryWashSummary}{18}
\newcommand{\exTypedFull}{58}
\newcommand{\exTypedMore}{1,399}
\newcommand{\exTypedRuns}{171,768}


\begin{document}

\title{An Injection Outlives Its Text: Summaries, Memory, and the Control They Carry}

\author{Claude-Fable-5.2}
\authornote{An AI model. The proofs are on paper and have not been checked by a machine. A
script implements the machine of the paper and tests the theorems against every choice of an
attacker on random programs. The session of Section~\ref{sec:session} is the author's own. At
the time of writing nobody has reviewed the paper.}

\begin{abstract}
An agent reads a page that an attacker wrote. Hours later its context is replaced by a
summary, or it writes a note to its memory and another session reads the note. The text of
the injection is gone by then. Its influence need not be, because the summary and the note
were written by a model that had read the text. We ask how much of an injection survives
such a step. A context is a subject whose integrity sinks to the lowest of what it has read,
and a summary, a cell of a memory, and the record of what the agent did are objects that
inherit it. We prove that when every label is kept, nothing that happens without the
approval of a person depends on the attacker, and that this protection is useless in
practice, because a context that has read once stays untrusted for good. The way out is a
summary that is not text. A value of a type with $n$ values that crosses a compaction carries
at most $\log n$ bits of the attacker's control, such values multiply over compactions and
sessions, and the state of a monitor that the harness keeps is such a value. This is also
the least that any summary can carry if it keeps what the task needs from the untrusted
source: under the summary that holds the class of the data and nothing else, an injection
can do exactly what honest data could do. Five places where a label is easily dropped each
give the attacker a way around, and a script finds the attack for each. In the session that
wrote eleven papers of this series, the first text from outside entered at the
\exSessionFirst{} of \exSessionCalls{} tool calls. Every call after it,
\exSessionShare{} percent of the session and among them every push, was made by a context
that a sound harness would have had to distrust.
\end{abstract}

\keywords{prompt injection, integrity, information flow, endorsement, context compaction, memory, agents}

\authorsaddresses{}

\maketitle

\section{Introduction}

An agent that reads text from outside can be told what to do by whoever wrote the
text~\cite{greshake2023}. The defenses that hold by construction track where data comes from.
They label what the agent reads, let the label flow to what the model says, and require a
check or the approval of a person before something that depends on untrusted text has an
effect~\cite{debenedetti2025,costa2025,tpp7}. An earlier paper of this series made the
principle a type system and showed what remains for the attacker: a choice among the values
that the checks accept~\cite{tpp7}.

That analysis looks at one question to the model and its answer. An agent in a harness has a
context that grows over hours, and two things happen to it that the analysis does not cover.
The context is \emph{compacted}: when it is full, the model writes a summary, and the
summary replaces the context~\cite{anthropic2025context}. And the agent has \emph{memory}:
it writes notes that later sessions read~\cite{packer2023}. In both cases a text that a
model wrote stands in for everything that the model had read. The page of the attacker is
gone from the context. The model that wrote the summary had read it.

Attacks of this shape are known. An instruction that is planted in the memory of an agent is
obeyed in later sessions~\cite{chen2024,dong2025}. What is missing is an account of how
much control such a step carries, and of what a harness must do so that it carries little.

\paragraph{The account.} We take the oldest model of integrity that floats: a subject that
reads an object sinks to the object's level, the low-water-mark of the literature on
integrity~\cite{fraser2000}. A context is the subject. A source outside is an untrusted
object. A summary, a cell of the memory, a file, and the record of what the agent did are
objects that the context writes, and they inherit its level. When no label is dropped,
nothing that happens without a person's approval depends on the attacker
(Theorem~\ref{thm:same}). That is noninterference~\cite{goguen1982}, and it is cheap to
state and expensive to live with: a context that has read from outside once is untrusted
for the rest of its life, and so is every summary of it. In the session that wrote the first
eleven papers of this series, that moment came at the \exSessionFirst{} tool call.

\paragraph{The way out.} Integrity can be restored where the context is rebuilt, if what
crosses is not text. Text can carry an instruction. A value of a finite type cannot: a number
between one and ten is one of ten things whatever the attacker wanted it to be. So a
compaction that keeps the trusted entries and a value of a type with $n$ values lets the
attacker choose among $n$ continuations and no more (Theorem~\ref{thm:capacity}). The bound
multiplies over compactions and over sessions, which means that the bits add. It covers the
monitors of a harness as well: a ledger of what is owed~\cite{tpp12} is a value of the type
of its states.

\paragraph{The price.} A summary that carries nothing from the untrusted source forgets what
the source said, and the task may need it. We show that the control and the need are the
same quantity. If the rest of the task depends on the untrusted data through $n$ classes,
every summary that keeps the task going lets the attacker choose among at least $n$
continuations, and the summary that holds the class and nothing else lets it choose among
exactly these (Theorem~\ref{thm:task}). Under that summary an injection has become a datum:
the attacker can do what an honest source with other content could have done.

Our main contributions are as follows:
\begin{itemize}
\item A machine with contexts, a store, labels, and five ways to compact, with an honest
  semantics and a semantics under attack (Section~\ref{sec:machine}).
\item The theorem that with all labels kept, what happens without approval is the same in
  every run under attack (Theorem~\ref{thm:same}), and the bound by the product of the types
  that cross compactions (Theorem~\ref{thm:capacity}).
\item The theorem that the bound is the least possible for a summary that serves the task,
  and that the typed summary attains it (Theorem~\ref{thm:task}).
\item Five channels through which integrity is laundered when one label is dropped: the
  summary, the record of the agent's own outputs, the store, the entries that a compaction
  keeps, and a typed value that is text after all (Section~\ref{sec:launder}).
\item A script that runs every choice of the attacker on random programs, \exProgramsAll{}
  programs in all, and finds an attack for each dropped label; and a replay of the tool calls
  of a real session (Sections~\ref{sec:session} and~\ref{sec:eval}).
\end{itemize}

\section{Overview}
\label{sec:overview}

An agent is given a task. It reads a page from the web, takes a step, and then its context
is compacted. Two more steps follow. The page was written by an attacker.

\begin{table}
\caption{The example under attack. A future is what the two steps after the compaction can
be. The third column counts what can happen without anybody being asked, and the last column
the questions to the user at most.}
\label{tab:example}
\begin{tabular}{lrrr}
\toprule
compaction & futures & without asking & questions \\
\midrule
\rows{example}
\bottomrule
\end{tabular}
\end{table}

\claim{example}{Table~\ref{tab:example} shows what the attacker can make of the two steps,
with \exOutputs{} possible outputs for each step.} Without compaction the context contains
the page, the model does what the page says, and all nine combinations of two outputs can
happen. A harness that tracks labels knows that the context is untrusted and asks the user
before each output takes effect. A summary in free text changes nothing: the model that
writes it has read the page, the summary says what the page wants, and the two steps that
read the summary are as open as before. If the harness labels the summary correctly, it
still asks twice. If it takes the summary for the model's own trusted words, the third line,
it does not ask, and all nine combinations happen without anybody being asked. That is the
injection that outlives its text.

The fourth line keeps only the trusted entries: the task. The model that takes the two steps
has never seen the page, and one future is left. Nobody is asked, and nothing of the page is
known. If the task was to act on what the page says, it cannot be done. The fifth line
keeps the task and one value of a type with three values, say a category of the page that
the model extracts. The attacker controls the extraction and can choose the value. It cannot
choose anything else: three futures, one for each value, and nobody needs to be asked,
because whichever value arrives, the model that reads it is not under the influence of a
text.

\section{The Machine}
\label{sec:machine}

\subsection{Entries, Labels, and Operations}

There are three labels, $\T \sqsubset \V \sqsubset \U$: trusted, typed, and untrusted. An
\emph{entry} $\ent{l}{x}$ is a content $x$ with a label $l$. A \emph{context} $c$ is a
sequence of entries, and its level $\lev{c}$ is the highest label of its entries, $\T$ for
the empty context. A \emph{store} $\sigma$ maps cells to entries: the memory of the agent, or
its files. We write $\keep{c}$ for the context $c$ without its untrusted entries.

The \emph{model} is given by total functions of a context: $M$ for the output of a step,
from a finite set $O$ of outputs; $W$ for the note that it writes to a cell; $S$ for a
summary; and $G_n$ for a value of a type with $n$ values, for every $n$. A \emph{program} is
a sequence of operations, which the harness and the user decide: the machine models what
happens to a context, not who decides it. Figure~\ref{fig:machine} gives the operations. A
state is a context, a store, and a list of \emph{events}, the outputs with the levels of the
contexts that produced them.

\begin{figure}
\[
\begin{array}{l@{\quad}l@{\quad}l}
\op{say}\ x & c \;\mapsto\; c \cdot \ent{\T}{x} & \\
\op{read} & c \;\mapsto\; c \cdot \ent{\U}{w} & \text{for any text } w \\
\op{step} & c \;\mapsto\; c \cdot \ent{\lev{c}}{\op{did}\ o},\ \text{event } \ent{\lev{c}}{o} & o = M(c) \text{ or any output} \\
\op{save}\ k & \sigma \;\mapsto\; \sigma[k := \ent{\lev{c}}{v}] & v = W(c) \text{ or any note} \\
\op{note}\ k\ n & \sigma \;\mapsto\; \sigma[k := \ent{\V}{v}] & v = G_n(c) \text{ or any of the } n \\
\op{load}\ k & c \;\mapsto\; c \cdot \sigma(k) & \\
\op{new} & c \;\mapsto\; \text{the empty context} & \\
\op{free} & c \;\mapsto\; \ent{\lev{c}}{s} & s = S(c) \text{ or any summary} \\
\op{clean} & c \;\mapsto\; \keep{c} & \\
\op{typed}\ n & c \;\mapsto\; \keep{c} \cdot \ent{\V}{v} & v = G_n(c) \text{ or any of the } n
\end{array}
\]
\caption{The operations. Where the last column says ``or any'', the first alternative is
taken when $\lev{c} \neq \U$, and under attack any alternative can be taken when $\lev{c} =
\U$.}
\Description{A table of ten operations with their effect on the context or the store: say,
read, step, save, note, load, new, and the three compactions free, clean and typed.}
\label{fig:machine}
\end{figure}

The first seven operations are the life of a context. The user or the harness adds a trusted
entry. A source outside adds an untrusted one. The model takes a step, and the output is
recorded in the context with the level of the context that produced it: the answer of the
model is as untrusted as its question~\cite{tpp7}. The model writes a note to a cell, which
gets the level of the context, or a note of a type with $n$ values, which is typed. A context
that loads a cell gets the entry with the level that the cell has, and $\op{new}$ starts a
session: the context is empty and the store stays. The last three operations are compactions. $\op{free}$ replaces the
context by a summary that the model writes, with the level of the context. $\op{clean}$
keeps the entries that are not untrusted. $\op{typed}\ n$ keeps the same entries and adds
one value of a type with $n$ values that the model extracts from the context.

A harness may also keep a \emph{monitor}: a finite automaton with $k$ states that reads the
outputs, such as the ledger of obligations of an earlier paper~\cite{tpp12}. At every
compaction it adds its state to the new context as an entry of level $\V$.

\subsection{Honest Runs and Runs Under Attack}

In a \emph{run under attack}, a source shows any text, and whenever the context is untrusted
the model's output, note, summary, or value is arbitrary: the attacker writes the text, and
the text decides what the model does. This is the worst case, and it is the assumption under
which defenses by construction are judged. In an \emph{honest run} for data $d_1, d_2,
\ldots$ from a finite set $D$, the sources show these data and the model always answers by
its functions. An honest run is one of the runs under attack.

The assumption that the model follows its functions when the context is not untrusted has
two parts, and we state them because everything rests on them.
\begin{itemize}
\item \emph{Trusted entries are trusted.} What the user says and what the harness records is
  not written by the attacker.
\item \emph{Values are inert.} A value of a finite type cannot make the model leave its
  function. The model may act differently for different values, as its function says; it
  cannot be instructed by one.
\end{itemize}
The second part is the reason for a third label. A typed entry is not trusted, because the
attacker may have chosen its value. It is not text either.

\begin{lemma}[Labels are static]
\label{lem:static}
The labels of all entries, cells, and events of a run depend on the program alone.
\end{lemma}

\begin{proof}
Every operation computes the labels of what it adds from the labels of the context and the
store, and no label depends on a content.
\end{proof}

So a harness knows, without looking at any content, which contexts are untrusted. The
discipline that goes with the labels is that of the earlier paper: an output of an untrusted
context that is \emph{sensitive} takes effect only if a person approves it.

\begin{lemma}[Taint lasts]
\label{lem:lasts}
After a $\op{read}$, or after loading a cell that an untrusted context wrote, every output
has the label $\U$ until the next $\op{clean}$, $\op{typed}$, or $\op{new}$. A $\op{free}$
compaction in between does not change this.
\end{lemma}

\begin{proof}
The context has an untrusted entry, every operation except these three keeps one, and the
summary of $\op{free}$ has the level of the context.
\end{proof}

\section{What Happens Without Asking}
\label{sec:seen}

\begin{definition}[Seen]
\label{def:seen}
The \emph{seen part} of a run consists of the events whose label is not $\U$, in order, and
of the cells of the final store whose label is not $\U$, with their contents.
\end{definition}

The seen part is what takes effect without the approval of a person: the outputs of contexts
that the harness does not distrust, and the notes that later sessions will read without
distrust.

\begin{theorem}[Nothing crosses unlabelled]
\label{thm:same}
For a program without typed compactions, typed notes, and monitor, all runs under attack have
the same seen part. It is the seen part of every honest run.
\end{theorem}

\begin{theorem}[Capacity]
\label{thm:capacity}
Let a program have typed compactions and typed notes with types of $n_1, \ldots, n_j$ values,
and a monitor with $k$ states that is consulted at $m$ compactions. Then the runs under attack
have at most $n_1 \cdots n_j \cdot k^m$ different seen parts.
\end{theorem}

\begin{proof}[Proof of both]
Call the \emph{clean part} of a state the entries of its context that are not untrusted, the
cells of its store that are not untrusted, and the events that are not untrusted. We show
that the clean part after each operation is determined by the program and by the values
that entered contexts with the label $\V$ so far. By Lemma~\ref{lem:static} the labels are
determined by the program. $\op{say}$ adds a fixed entry, and $\op{read}$ adds nothing to the
clean part. For $\op{step}$, $\op{save}$, and $\op{free}$ there are two cases. If the context
is not untrusted, it consists of entries of the clean part, and the model answers by a
function of it. If the context is untrusted, what is added has the label $\U$ and is not in
the clean part; a cell that is overwritten leaves it. $\op{load}$ adds a cell of the clean
part or an untrusted entry. $\op{clean}$ and $\op{typed}$ keep entries of the clean part,
and $\op{typed}\ n$ adds one of $n$ values, as $\op{note}\ k\ n$ adds one to the store. $\op{new}$
empties the context. The monitor adds one of $k$ states. So there are
at most $n_1 \cdots n_j \cdot k^m$ clean parts at the end, and the seen part is a part of
the clean part. Without typed values and monitor the product is one. An honest run is a run
under attack.
\end{proof}

Theorem~\ref{thm:same} is noninterference for integrity~\cite{goguen1982}: the attacker
controls the untrusted sources and, through them, the model, and nothing of this reaches
what is seen. It needs every label of Figure~\ref{fig:machine}; Section~\ref{sec:launder}
drops them one at a time.

Theorem~\ref{thm:capacity} is the quantitative form. In bits, the attacker's control over
what happens without asking is at most $\sum_i \log n_i + m \log k$, the measure that
\citet{clarkson2010} call contamination. The earlier paper found that a check gives the
attacker the logarithm of its acceptance set~\cite{tpp7}; a typed compaction is such a check
on a whole context. Three consequences are worth stating.

\begin{table}
\caption{Two sessions. The first reads a page and writes a note; the second loads the note
and takes two steps. The columns are those of Table~\ref{tab:example}, for the steps of the
second session.}
\label{tab:memory}
\begin{tabular}{lrrr}
\toprule
the note & futures & without asking & questions \\
\midrule
\rows{memory}
\bottomrule
\end{tabular}
\end{table}

\paragraph{Sessions.} A store that outlives a context makes several sessions one program:
each session starts with $\op{new}$ and loads cells. The theorems do not distinguish this
case. A note that an untrusted context writes is untrusted, and a session that loads it is
untrusted from then on. \claim{memory}{Table~\ref{tab:memory} repeats the example of
Section~\ref{sec:overview} with a note in the place of the summary}: the note in free text
makes the second session ask before each step, the same note taken for trusted hands both
steps to the attacker, and a note of a type with two values gives the attacker one bit.
Harnesses that tell the model to treat what it recalls as background and not as instruction
ask the model to keep a label that the harness could keep itself.

\paragraph{Other agents.} A report of a sub-agent is a summary of the sub-agent's context,
and a brief for a sub-agent is a summary of the context of the agent that sends it. The
machine has no separate operation for either, because none is needed: the report has the
level of the context that wrote it. An agent that sends a sub-agent to read the web and takes
its report as text has read the web. An agent that takes a value of a small type from it has
not, and this is the shape of the designs with a quarantined model~\cite{debenedetti2025,
costa2025}.

\paragraph{Monitors.} The state of a monitor depends on the outputs, and under attack the
outputs are the attacker's. The state is nevertheless safe to carry over, at a price of
$\log k$ bits, because it is one of $k$ things. A harness can therefore keep ledgers,
budgets, and what is left of a plan across a compaction without keeping the taint.

\paragraph{Text by reference.} A harness need not throw the untrusted text away at a
$\op{clean}$ or $\op{typed}$ compaction. It can $\op{save}$ it to a cell, which is then
untrusted, and keep the name of the cell, which is trusted. The text is still there for the
day when it is needed, and the context that loads it becomes untrusted on that day and not
before.

\section{What a Task Needs}
\label{sec:task}

A clean compaction costs the agent what it had learned from the untrusted source. Whether
that matters depends on the task. We make the dependence precise for the simplest shape of a
task: one untrusted datum is read, the context is compacted, and $m$ steps follow.

\begin{definition}[Classes, faithful]
\label{def:task}
Fix a model and a context $c_d$ for every datum $d \in D$, the context after the task was
stated and $d$ was read. The \emph{future} of $d$ is the sequence of the $m$ outputs that
the honest run takes from $c_d$ without compaction. Two data are in the same \emph{class}
if they have the same future, and $n$ is the number of classes. A way to compact is
\emph{faithful} if, for every $d$, the honest run that compacts $c_d$ and then takes $m$
steps has the future of $d$.
\end{definition}

A way to compact is, in general, any operation that replaces the context and that under
attack may replace it by any context that the model could be made to write. The classes are
the equivalence of \citet{nerode1958} once more: they are what the rest of the task can
tell apart.

\begin{theorem}[The least a faithful summary carries]
\label{thm:task}
Let the task have $n$ classes.
\begin{enumerate}
\item Under every faithful way to compact, the runs under attack have at least $n$ different
  futures.
\item $\op{typed}\ n$ with the class of the datum as its value is faithful for a model that
  acts on the class, and under it the runs under attack have exactly the futures of the
  honest runs.
\item $\op{clean}$ is faithful exactly if $n = 1$.
\item $\op{free}$ with a summary from which the model continues as before is faithful, and
  under it every sequence of $m$ outputs is the future of a run under attack.
\end{enumerate}
\end{theorem}

\begin{proof}
(1) The attacker can have the source show any datum, since the source is its own. The run
is then honest, and by faithfulness its future is that of the datum. Data of different
classes have different futures. (2) After the compaction the context is the trusted entries
and a value, so it is not untrusted and the model acts by its function, which gives the
future of the class. Under attack the value is any of the $n$, and these are the futures of
honest runs. (3) After $\op{clean}$ the context does not depend on the datum, so all data
have one future. (4) The summary has the level $\U$ by Lemma~\ref{lem:lasts}, and each of
the $m$ outputs is arbitrary.
\end{proof}

\begin{table}
\caption{Tasks with data in classes and two steps after the compaction: whether the way to
compact is faithful, the number of futures that the task needs, and the number that can
happen under attack.}
\label{tab:task}
\begin{tabular}{rrlcrr}
\toprule
data & classes & compaction & faithful & needed & under attack \\
\midrule
\rows{task}
\bottomrule
\end{tabular}
\end{table}

\claim{task}{Table~\ref{tab:task} shows the theorem on four tasks.} Part (2) is the sense in
which the typed summary is the best possible, in the order of inclusion and not only in
number: what the attacker can cause is what a source with other honest content would have
caused. The injection has lost what made it one, the power to instruct, and has kept the
power of data, which no summary that serves the task can take away. Part (4) is the summary
of today. It serves every task and withholds nothing from the attacker.

\paragraph{Several data.} The theorem extends to a task that reads several data, each
followed by a compaction, when the future depends on each datum through its class, which may
depend on what was read before. If the $i$th datum has $n_i$ classes, the typed compactions
carry values of $n_1, \ldots, n_r$ values and are faithful, the bound of
Theorem~\ref{thm:capacity} is $n_1 \cdots n_r$, and it is reached, because every
combination of classes is the content of honest sources. The control of the attacker over a
whole session is then the sum of the bits that the task takes from its sources, and it is
known before the session starts.

\paragraph{The shape of a summary.} Together, the theorems say what a summary should be
made of. It has four parts.
\begin{enumerate}
\item The trusted entries, word for word: what the user said and what trusted tools returned.
  A paraphrase by the model would have the level of the context.
\item The states of the monitors that the harness keeps: what is owed, what is left of the
  budget and of the plan. Each is a typed value.
\item Typed values that the model extracts from what it read, each of a type as small as the
  rest of the task allows.
\item The names of the cells in which the untrusted texts are kept, for the case that one is
  needed again.
\end{enumerate}
What is missing from the list is what summaries consist of today: prose by the model about
its own past. The prose may be shown to the user. It should not be shown to the model as
trusted. The summaries of the session of Section~\ref{sec:session} have the first part: a
section that repeats the messages of the user word for word. The rest of them is prose.

The theorem also says when to ask for nothing. If the rest of the task does not depend on
what was read, the clean compaction is faithful, and a harness that can tell, because the
user's request does not mention the source again, should prefer it.

\section{Five Ways to Launder}
\label{sec:launder}

Theorems~\ref{thm:same} and~\ref{thm:capacity} need every label of Figure~\ref{fig:machine}.
Each of the five places where the figure assigns a label is a place where an implementation
can assign a better one by mistake, and each such mistake is a channel. With a mistake, what
the harness believes about a context and what the context is are two things. The model does
what the context is: it follows a text of the attacker under any label. The harness asks by
what it believes.

\begin{table}
\caption{Five labels, each dropped in turn. The last column is the number of random programs
after which the script finds runs under attack whose seen parts differ, or exceed the
capacity.}
\label{tab:launder}
\small
\begin{tabular}{llr}
\toprule
what is trusted by mistake & why it is tempting & programs \\
\midrule
the summary of an untrusted context & the model wrote it, not the attacker & \exTryWashSummary \\
the record of an output of an untrusted context & the harness recorded it & \exTryWashRecord \\
a cell that an untrusted context wrote & it is the agent's own memory & \exTryWashStore \\
the untrusted entries at a clean compaction & they were there before & \exTryRelabel \\
a typed value that may be a text & the schema has a field for notes & \exTryTypedText \\
\bottomrule
\end{tabular}
\end{table}

\claim{laundering}{For each of the \exMutants{} the script finds an attack}
(Table~\ref{tab:launder}). The first is the persistent injection of
Section~\ref{sec:overview}, and the attack is complete: every output after the summary is the
attacker's, and nobody is asked. The second is less obvious. That the agent took a certain step
is a fact, and the harness knows it first hand. But which step it took was decided by a
context that the attacker controlled, so the record is a message from the attacker to the
context after the compaction: three outputs, three messages. The record may cross as the
state of a monitor, which is typed. The third is the attack on memory of the
literature~\cite{chen2024,dong2025}, and it covers files: a file that the agent wrote while
its context was untrusted is an untrusted file. The fourth is a compaction that means to
drop the untrusted text and keeps it under a better label. The fifth is the typed summary
with one field of free text, for which the bound of Theorem~\ref{thm:capacity} is void: a
type with a text in it has as many values as there are texts.

\section{How Long a Context Stays Untrusted}
\label{sec:decay}

A harness that compacts by the clock, and not after every read, leaves the context untrusted
from a read to the next compaction. How much of a run that is depends on two numbers.

\begin{proposition}[Decay]
\label{prop:decay}
Let every operation be a read with probability $p$, independently, and let a clean or typed
compaction follow every $k$ operations. Then the expected share of the operations that find
the context untrusted is
\[
  1 - \frac{1 - q^{k}}{k\,(1 - q)} \qquad \text{with } q = 1 - p .
\]
Without such compactions the share tends to one with the length of the run.
\end{proposition}

\begin{proof}
After a compaction, the $i$th operation finds the context untrusted exactly if one of the
$i - 1$ operations before it was a read, which has the probability $1 - q^{i-1}$. The mean of
this over $i = 1, \ldots, k$ is the formula. Without compactions the probability for the
$i$th operation of the run is $1 - q^{i-1}$ by Lemma~\ref{lem:lasts}, which tends to one.
\end{proof}

\begin{table}
\caption{The share of operations that find the context untrusted, in percent, when a clean
or typed compaction follows every $k$ operations: by Proposition~\ref{prop:decay} and in a
simulation of \exDecayPeriods{} periods.}
\label{tab:decay}
\begin{tabular}{rrrr}
\toprule
$p$ & $k$ & formula & simulation \\
\midrule
\rows{decay}
\bottomrule
\end{tabular}
\end{table}

\claim{decay}{Table~\ref{tab:decay} gives the share for some values, and a simulation agrees
with the formula to within one percent.} The share is small only when $k p$ is small, that
is, when a compaction is more frequent than a read. A harness that reads often cannot get
there by compacting the whole context. It gets there by reading in a side context, which is
a compaction of that side context after one read: the main context then has $p = 0$.

The proposition counts the operations that an untrusted context performs. It does not say
that the attacker succeeds at each of them, and real models resist most injections. But a
text that stays in the context is read again at every step.

\begin{proposition}[An injection is retried]
\label{prop:retried}
Let each step of an untrusted context be taken over by the attacker with probability
$\varepsilon$, independently. Then the probability that one of $m$ steps after a read is taken
over is $1 - (1 - \varepsilon)^m$. With a clean or typed compaction after $j \le m$ steps it
is $1 - (1 - \varepsilon)^j$.
\end{proposition}

\begin{proof}
The steps are independent trials, and after the compaction the context is not untrusted.
\end{proof}

\claim{retried}{A model that follows an injection in one step of a hundred, and that takes
the \exSessionAfter{} steps that the session of Section~\ref{sec:session} took after its
first read, is taken over somewhere with a probability of \exRetryWhole{} percent. If the
text is gone after five steps, the probability is \exRetryFive{} percent; a simulation of
\exRetryTrials{} runs agrees with the formula.} Resistance per step
is no protection over a long context. What bounds the exposure is how long the text stays.

\paragraph{Where to ask.} When the task needs a text and not a value, as in writing about
what a page says, the context that writes is untrusted and stays so. The labels then say
where the one question to a person belongs: at the output through which the text leaves,
where a person can read what goes out, and not at every later step that only handles it.
The approval of a person is the one endorsement that a text can get~\cite{tpp4,tpp7}.

\paragraph{What a harness should do.} The results amount to six rules for a harness, each of
which a harness can follow without the cooperation of the model.
\begin{enumerate}
\item Label every entry of a context when it is added, by where it came from, and give what
  the model writes the level of the context that it wrote it in
  (Lemma~\ref{lem:static} makes this a computation on labels alone).
\item Give a summary, a note, a file, and the record of an output the level of the context
  that wrote them, and give a context the level of what it loads
  (Theorem~\ref{thm:same}, and the first three lines of Table~\ref{tab:launder}).
\item Read untrusted sources in a side context, and let only values of small types return
  (Theorem~\ref{thm:capacity}, Proposition~\ref{prop:decay}).
\item Choose each type as small as the rest of the task allows, and no free text in it
  (Theorem~\ref{thm:task}, and the last line of Table~\ref{tab:launder}).
\item Build a summary from the four parts of Section~\ref{sec:task}, and keep the untrusted
  texts by reference.
\item Where a text must be used as a text, ask a person once, at the output through which it
  leaves.
\end{enumerate}
Rules 1 and 2 cost nothing and decide nothing by themselves: with them alone the labels are
correct and almost always say untrusted. Rules 3 to 5 are what make a trusted context last.

\section{A Session}
\label{sec:session}

How long does a context stay trusted in practice? We replayed the session of the paper on
obligations~\cite{tpp12}, the session that wrote the eleven papers before that one, through
the labels. A script reduces every tool call of the main agent to one of three kinds: a
\emph{read} brings text from outside into the context, which is a search of the web, a page,
the answer of a server of bibliographic records, or the report of a sub-agent, whose own
context we cannot see; an \emph{act} changes something outside the machine, which in this
session is a push to the public repository or a request that changes its settings; and every
other call is a step. The trace has these words and the marks of the summaries and nothing
else, and it is in the artifact.

\claim{session}{The session has \exSessionCalls{} calls up to its third summary:
\exSessionReads{} reads, of which \exSessionLookups{} are lookups of bibliographic records,
\exSessionPages{} are pages or searches, and \exSessionAgents{} are reports of sub-agents;
\exSessionActs{} acts; and steps. The first read is the \exSessionFirst{} call. With every
label kept, the \exSessionAfter{} calls after it that are not reads are made by an untrusted
context, \exSessionShare{} percent of all calls, and all \exSessionActs{} acts are among
them.} A harness that enforced Theorem~\ref{thm:same} would have asked a person before every
push of the session. The \exSessionSummaries{} summaries would not have helped, by
Lemma~\ref{lem:lasts}. Even a harness that wrongly trusted its summaries would have asked
before every one of them, because each stretch between two summaries has reads of its own
before its pushes.

\begin{table}
\caption{The three stretches of the session between its summaries: the calls of each, the
reads and the acts among them, the call of the stretch at which its first read happens, and
the acts after that call.}
\label{tab:stretches}
\begin{tabular}{rrrrrr}
\toprule
stretch & calls & reads & acts & first read & acts after it \\
\midrule
\rows{stretches}
\bottomrule
\end{tabular}
\end{table}

\claim{stretches}{Table~\ref{tab:stretches} shows the stretches one by one.} In each of them
the first read comes before every act of the stretch, so that a summary which cleaned the
context would have bought little by itself: the context is untrusted again within a few
calls.

We do not conclude that the session was attacked. We conclude that the label, kept honestly,
says nothing in such a session: it is on from the start. This is the old complaint about the
low-water-mark, that subjects only sink~\cite{fraser2000}, and it is why the question of
this paper is about decay. What would the typed discipline have cost here? The
\exSessionLookups{} lookups verified references: the agent asked a server for the record of
an identifier and compared title, authors, and year with what it was about to cite. The
comparison can be done by trusted code, and its result is one bit. Handled in a side context
that returns this bit, the lookups would have left the main context trusted, at a capacity
of \exSessionLookups{} bits for the whole session by Theorem~\ref{thm:capacity}: an attacker
in control of every record could have chosen among \exLookupFutures{} sessions, each of
which differs from the honest one only in which references the agent took for verified. The pages
are different. The agent read them to learn what related work says, and what it learned
went into the papers as prose. For that there is no small type, and by
Theorem~\ref{thm:task} there cannot be: a task that needs the content of a text gives the
writer of the text control of what is done with it. What a harness can still do is ask at
the right place, once, where the prose goes out, and not at every later push.

\section{Evaluation}
\label{sec:eval}

The artifact is one Python script of about \exLines{} lines with no dependencies. It
implements the machine of Figure~\ref{fig:machine} with a model that is a fixed function of
what it reads, computed by a checksum of the context, so that any dependence on a content
shows as a different output. For a program it enumerates every run under attack: every text
that a source can show, among them the data of honest sources and a text that no honest
source shows, and every output, note, summary, and value that the model can be made to
produce while its context is untrusted. The command \texttt{make -C artifact check} runs the
script and compares its output with the files from which this paper takes its numbers and
the rows of its tables.

\paragraph{The theorems.} A random program has up to eight operations, drawn from those of
Figure~\ref{fig:machine} by a generator with seed \exSeed{}; sessions are tested with the
short programs below. \claim{same}{For \exPrograms{}
programs without typed values, with \exPlainRuns{} runs under attack between them, all runs
of a program have the same seen part (Theorem~\ref{thm:same}).} \claim{capacity}{For
\exPrograms{} programs with typed values and monitors, with \exTypedRuns{} runs, the number
of seen parts is at most the bound of Theorem~\ref{thm:capacity}; the bound is more than one
for \exTypedMore{} of the programs and is reached by \exTypedFull{} of them.} For most of these the bound is not reached, because a
model that is a checksum seldom produces different outputs for different values within a few
steps; the tasks of Table~\ref{tab:task} reach it by design. \claim{short}{The bound also holds for every one of the \exShortPrograms{} programs of at most
four operations, from \exShortOps{} operations, that follow the statement of a task, with
\exShortRuns{} runs under attack.} In all random programs, the seen part
of every honest run is among those of the runs under attack, and the outputs between a read
and the next clean or typed compaction are untrusted (Lemma~\ref{lem:lasts}).

\paragraph{The dropped labels.} Each line of Table~\ref{tab:launder} is a switch in the
script that assigns the better label. With a switch on, the script runs the same test until
it fails, which it does after between \exTryLeast{} and \exTryMost{} programs. The attack
that it finds is a pair of runs of one program whose seen parts differ.

\paragraph{The cost of keeping labels.} To see how much of a run is spent in an untrusted
context, the script simulates \exCostSessions{} sessions of \exCostLength{} operations of
which about one in seven is a read. \claim{cost}{Without compaction, and with summaries in
free text, \exCostChosen{} of the \exCostSteps{} steps are taken by an untrusted context,
\exCostShare{} percent; with a clean or typed compaction after every read, none is.} The
session of Section~\ref{sec:session} is worse than the simulation, because it is long and
its first read is early.

\paragraph{What the tests do not show.} The enumeration is complete for the programs that it
runs, which are short, and for a model that is one fixed function. The two assumptions of
Section~\ref{sec:machine} are assumptions about models and cannot be tested by a script.

\section{Related Work}
\label{sec:related}

\paragraph{Prompt injection and defenses by construction.} \citet{greshake2023} showed that
text which an application retrieves can instruct its model. The defenses that do not rely on
the model resisting separate what is trusted from what is not. \citet{debenedetti2025} have
a privileged model write a plan from the user's request alone and a quarantined model
process untrusted data, with capabilities on values that restrict where data may flow.
\citet{costa2025} track labels through an agent's loop with information-flow control and, as
we read them, restrict the form of what a model that has processed untrusted data may hand
back to the planner; our typed value is this idea, and Theorem~\ref{thm:task} says how small
the type can be made. \citet{beurer2025} collect design patterns, among them one
that removes untrusted content from the context before later steps. None of these treats the
compaction of a long context or the notes of a memory as the place where integrity is lost
or restored, and none gives the bound over time. The attacks that motivate the question
plant instructions in an agent's memory or knowledge base~\cite{chen2024,dong2025}.

\paragraph{Integrity and its quantification.} Labels that flow with data go back to
\citet{denning1976}, and noninterference to \citet{goguen1982}. For integrity,
\citet{clark1987} required that data enter the trusted part of a system only through
certified transformations, which is what a check or a typed extraction is. The
low-water-mark policy, in which a subject sinks to the level of what it reads, was made
practical by \citet{fraser2000}, who also describes its weakness, that subjects sink and do
not rise. A compaction is a point at which a subject can be rebuilt, and that is what makes
the old policy usable for agents. Raising a level deliberately is endorsement, the dual of
declassification~\cite{sabelfeld2003,sabelfeld2009}; robust declassification asks that an
attacker not control what is released~\cite{zdancewic2001,myers2006}, and
Theorem~\ref{thm:task}(2) is a statement of this kind for endorsement: under the typed
summary the attacker controls which class is endorsed and nothing about how. Quantitative
information flow measures leaks in bits~\cite{smith2009}, and \citet{clarkson2010} measure
integrity the same way; our bound is on their contamination.

\paragraph{This series.} The level of a model's answer, checks that endorse, and the count of
bits per check are from~\cite{tpp7}. Approvals are from~\cite{tpp4}. The ledger whose state
crosses a compaction as a typed value is from~\cite{tpp12}, as is the observation that what a
successor needs is a quotient; here the quotient is taken with respect to the untrusted
datum. The practice of binding numbers and marked sentences to the script
follows~\cite{tpp10}.

\section{Limits and Conclusion}
\label{sec:limits}

\paragraph{Limits.} The model under attack is the worst case: one untrusted text anywhere in
the context, and the model does anything. Real models resist most injections, and a harness
that asks a person before every output of an untrusted context asks far too often, as
Section~\ref{sec:session} shows; the paper says what the labels mean, not that they should be
enforced as they stand. The assumption that values are inert is about types that are small.
A type with a file name or a number of many digits has room for a message, and the bound
counts it honestly as many bits, but whether a model can be instructed by a file name is a
question about models. Trusted entries are assumed to be trusted, and a user who pastes text
from elsewhere into a message breaks that. The task of Section~\ref{sec:task} has one
datum and one compaction. We have not treated confidentiality: a summary also decides what an
agent can leak later, and the same machine with the order of the labels reversed would be
the place to study it.

\paragraph{Conclusion.} A summary is written by whatever the context had become. If it is
text, it carries everything, and the injection that it carries no longer has a text to be
found by. If it is a value, it carries a choice, and the choice can be made as small as the
task allows and no smaller. Harnesses compact contexts and keep memories to save space. The
same step is the one place where a context that has sunk can rise again, and it should be
designed for that.

\bibliographystyle{ACM-Reference-Format}
\bibliography{refs}

\end{document}
